\chapter{Introduction and model scope}\label{ch:introduction}
\section{What AMPS solves}
AMPS is a kinetic simulation framework for rarefied gas, neutral and charged
particle populations on adaptive meshes. Applications use its particle
records, transport, stochastic reactions/collisions, surface interactions,
sampling and distributed data structures. Depending on configuration, fields
are prescribed, coupled from another model, or advanced by supported PIC
solvers. The same framework supports planetary exospheres, comet comae and
dust, magnetospheric particles and heliospheric transport.

AMPS can run standalone or as an SWMF component. A standalone application
owns its outer lifecycle; SWMF supplies component communicators, scheduling
and coupling epochs through wrappers. Shared numerical machinery does not
make the scientific assumptions of one application valid in another.

\section{Choose the physical problem first}
A useful run description states the species, phase-space representation,
geometry/frame, initial and external boundary conditions, injection rate and
velocity law, force model, reaction network, surface physics and observable.
Then state mesh, time-step, particle-weight and sampling policies. If any
required physical information is absent, resolve it before running: numerical
convenience is not a physical closure.

For example, a neutral Enceladus parent-loss calculation can have H2O and CO2
particles while allocating no photoproducts and using no plasma. A comet
network can instead produce daughters with recoil and internal energy. A
field-line SEP model and a full 3-D charged-particle mover have different
transport approximations. A source-free analytical CME background test cannot
qualify particle acceleration or a full downstream plasma volume.

\section{Four kinds of evidence}
\begin{enumerate}
\item Structural checks verify grammar, registration, data identity and
supported configuration. They need not evolve any physical state.
\item Numerical verification compares an algorithm with independent equations,
manufactured solutions, conservation or convergence.
\item Regression testing detects change relative to reviewed fixtures and can
include a full configure/build/run/output path.
\item Scientific validation compares the declared physical model with relevant
laboratory, observational or independently justified physical reference data.
\end{enumerate}
These layers complement each other. A model can conserve particles while
using the wrong source law; it can reproduce a fitted observation without
independently validating that source. Reports must preserve the distinction.

\section{Reading and operating conventions}
Bare dimensional values in modern application runtime keys generally use the
SI units stated by the owner; older generators and external tables can have
explicit different units. The applicable input parser is authoritative.
A path shown in a command is relative to the repository root unless a working
directory is declared. Placeholders must be supplied by the operator.
The documented test commands use a required machine profile rather than
hard-coded compiler/kernel paths. Native build examples require an idle
checkout and a clean generated build as explained next.
\source{src/README.md; src/pic/README.md; src/test/README.md}
